\documentclass[a4paper,11pt,dvipdfmx]{jsarticle}
\usepackage{amsmath, amssymb, amsthm}
\usepackage{bm}
\usepackage{mathtools}
\usepackage{comment}
\usepackage{algorithm}
\usepackage{algorithmic}
\usepackage{bm}
\usepackage{mathtools}
\usepackage{enumitem}
\mathtoolsset{showonlyrefs}
\usepackage[dvipdfmx, colorlinks=true, linkcolor=blue, citecolor=blue, urlcolor=blue, setpagesize=false]{hyperref}
\usepackage{pxjahyper}
% --- 定理・定義環境 ---
\theoremstyle{definition}
\newtheorem{definition}{定義}[section]
\newtheorem{theorem}[definition]{定理}
\newtheorem{lemma}[definition]{補題}
\newtheorem{assumption}{仮定}
\newtheorem{proposition}[definition]{命題}
\newtheorem{requirement}[definition]{要請}
\newtheorem{example}[definition]{例}
\newcommand{\fU}{\underline{f}}
\newcommand{\gU}{\underline{g}}
\newcommand{\uU}{\underline{u}}
\newcommand{\aU}{\underline{a}}
\newcommand{\bU}{\underline{b}}
\newcommand{\cU}{\underline{c}}
\newcommand{\dU}{\underline{d}}
\newcommand{\vU}{\underline{v}}
\newcommand{\zeroU}{\underline{0}}
\newcommand{\C}[2]{C_{{#1},{#2}}}
\newcommand{\cy}[3]{\mathcal{C}_{{#1},{#2},{#3}}}
\newcommand{\inv}{^{-1}}
\newcommand{\HHX}{\hat{H}_X}
\newcommand{\HHZ}{\hat{H}_Z}
\newcommand{\tr}{^\top}

% =========================================
\begin{document}
\section{ブロックサイズ $P$ 縮小に向けた構成手法}

本章では、提案されたアフィン置換行列（APM）に基づく量子LDPC符号の構成において、ブロックサイズ $P$ を768未満に縮小するための2つのアプローチについて論じる。原論文の構成例においては $P=768$ が採用されているが、これは特定の可換性制約と内周（girth）条件を同時に満たすための代数的な制約に起因する。以下に、この壁を突破するための数学的および構造的な制約緩和のアプローチを提示する。

\subsection{アクティブインデックス集合の分散配置による制約緩和}

第一の手法は、アクティブ行列を構成するインデックス集合 $S$ の選択を変更し、可換性制約そのものを削減するアプローチである。

原論文の標準的な構成においては、アクティブな部分行列として親行列の上位3ブロック行（すなわち $S = \{0, 1, 2\}$）を選択している。この連続した配置を採用した場合、部分行列間の直交性を保証するために可換性が要求されるインデックス差分の集合 $\Delta$ のサイズは、 $|\Delta| = 5$ となる。可換性を満たすべき条件の多さが、結果として $P$ のサイズを押し上げる主要な要因となっている。

これに対し、論文の6.2節のRemarkにおいて示唆されている通り、アクティブな行の選択を $S = \{0, 2, 4\}$ のように分散して配置することが極めて有効である。この分散配置を採用した場合、可換性が要求される差分集合のサイズは $|\Delta_S| = 3$ にまで大幅に削減される。満たすべき代数的な制約条件（二次合同式）が減ることで、アフィン置換行列のパラメータ探索空間が緩和される。これにより、結果として $P < 768$ のより小さな次元においても、アクティブな直交性を維持したまま、条件を満たす $F, G$ のパラメータ集合を発見できる可能性が飛躍的に高まる。

\subsection{可換性テーブルの再設計と非可換ペアの拡張}

第二の手法は、探索の前提条件となっている「可換性テーブル」のパターン自体を再設計し、探索空間を拡張するアプローチである。

著者らは、BP復号において有害となるトラッピングセット（ETS）の生成を抑制しつつ、低重みの論理演算子がアクティブ行列の双対空間から漏れ出すのを防ぐために、特定の2つのペア（例えば $(0, 3)$ と $(1, 2)$）のみを非可換とするヒューリスティックな設計を採用している。

しかしながら、この特定のパターンに限定することは、探索空間を過度に狭める結果を招いている。非可換とするペアの組み合わせをさらに拡張し、3つ以上のペアを非可換とする、あるいは全く異なるインデックスの組み合わせを採用することにより、より小さな $P$ の次元において構成可能な新たなパラメータ空間が開拓される。この際、短いサイクル（girthの低下）や新たなトラッピングセットの発生リスクが生じるが、これらを探索プログラムの目的関数に組み込んだ上で探索を行うことで、要件を満たすよりコンパクトな行列構成の発見が期待される。

\section{アクティブインデックス集合の分散配置に基づく探索手順の具体化}

前章で提示した「アクティブインデックス集合 $S$ の分散配置」について、探索アルゴリズムに組み込むための具体的な手順を定式化する。本手法の核心は、抽出するブロック行のインデックスを意図的に離散させることで、直交性を維持するために必要な可換性条件（インデックス差分集合 $\Delta$ のサイズ）を最小化することにある。

\subsection{インデックス集合と差分集合の再定義}
原論文の設定である列重み $J=3$、親行列のハーフサイズ $L/2=6$ を前提とする。標準的な構成では、アクティブインデックス集合は上位3行 $S = \{0, 1, 2\}$ であり、要求される差分集合は以下の通りであった。
\[
\Delta = \{(k - i) \pmod 6 \mid i, k \in \{0, 1, 2\}\} = \{0, 1, 2, 4, 5\}
\]
この制約を緩和するため、インデックス集合を等間隔に分散配置した $S' = \{0, 2, 4\}$ に再定義する。この変更により、新たな差分集合 $\Delta_{S'}$ は次のように算出される。
\[
\Delta_{S'} = \{(k - i) \pmod 6 \mid i, k \in \{0, 2, 4\}\} = \{0, 2, 4\}
\]
結果として、可換性が必須となる条件の数が5個から3個へ減少し、探索空間の自由度が大幅に向上する。

\subsection{アルゴリズム実装における具体的手順}
探索プログラムにおいて、この理論的背景を実装に落とし込む手順は以下の4段階で構成される。

\begin{enumerate}
    \item \textbf{アクティブ行の抽出ロジックの変更}\\
    親行列 $\hat{H}_X$ および $\hat{H}_Z$ からアクティブ行列 $H_X, H_Z$ を構築する際、取得するブロック行のインデックスを $0, 1, 2$ から $0, 2, 4$ へ変更する。これにより、Tannerグラフを構成するノードの接続関係（パリティ検査行列の構造）が直接的に再構成される。
    
    \item \textbf{動的制約伝播における判定条件の緩和}\\
    変数候補を生成する際、線形合同式を解いて探索空間を絞り込む基準となる $\Psi_r = 0$ の判定条件を、$r \in \{0, 1, 2, 4, 5\}$ から $r \in \{0, 2, 4\}$ に縮小する。この緩和により、各ステップで有効となる候補ペア $(a, b)$ の数が大幅に増加する。

    \item \textbf{非可換パラメータの許容範囲の拡張}\\
    親行列の直交性を破り、潜在的な低重み論理演算子を回避するための非可換条件 $\Psi_r \neq 0$ を設定できるインデックス $r$ の候補が、従来の $\{3\}$ のみから、$\{1, 3, 5\}$ の3種類へと拡張される。探索時には、これら複数の非可換成分を組み合わせた多様な可換性テーブルを生成・評価する。

    \item \textbf{サイクル検出およびトラッピングセット評価の再計算}\\
    インデックス集合の変更に伴い、アクティブ行列内に形成される短いサイクル（8サイクルなど）の発生パターンや、有害なトラッピングセットの構造が変化する。したがって、探索アルゴリズム内の \texttt{check\_cycles} 関数やトラッピングセットの事前計算ライブラリを、新たなインデックス構成 $S' = \{0, 2, 4\}$ に基づいて完全に再評価し、枝刈りの基準を更新する。
\end{enumerate}

これらの手順を適用することで、不必要な代数的制約による探索空間の枯渇を防ぎ、$P=768$ 未満のより小さなブロックサイズにおいても、指定された内周条件と直交性を満たす解への到達が期待できる。

\subsection{アクティブ部の直交性と潜在部の非直交性の両立}

インデックス集合を分散配置し $S' = \{0, 2, 4\}$ とした場合、探索空間を拡大するだけでなく、量子LDPC符号として最低限満たすべき「アクティブ行列間の直交性」と、低重み論理演算子を防ぐための「潜在行列とアクティブ行列間の非直交性」を同時に成立させなければならない。以下に、これら二つの条件を両立するための具体的な手順と数理的条件を定式化する。

\subsubsection{アクティブ部の直交性条件}
アクティブ行列 $H_X$ および $H_Z$ は、親行列の行インデックス $S' = \{0, 2, 4\}$ を抽出して構成される。CSS符号のパリティ検査行列として機能するためには、アクティブ直交性 $H_X H_Z^T = 0$ が必須である。
ブロック行列の積の性質より、抽出されたブロック $(i, k) \in S' \times S'$ における相互作用行列 $\Psi_r$ は、そのインデックス差分 $r = (k - i) \pmod 6$ に依存する。
$S'$ 内の要素同士の差分集合は、
\[
\Delta_{S'} = \{(k - i) \pmod 6 \mid i, k \in \{0, 2, 4\}\} = \{0, 2, 4\}
\]
となる。したがって、アクティブ部の直交性を保証するための必要十分条件は、以下の3つの相互作用行列が零行列となることである。
\[
\Psi_0 = 0, \quad \Psi_2 = 0, \quad \Psi_4 = 0
\]
探索アルゴリズムにおいては、差分が $0, 2, 4$ となるすべてのペア $(F_u, G_{r-u})$ に対して可換性を強制する。

\subsubsection{潜在部の非直交性条件と混合ブロックの構造}
次に、潜在行列 $\tilde{H}_X, \tilde{H}_Z$ がアクティブ行列と直交してしまうこと（親行列全体の直交性）を防ぐ条件を定式化する。潜在行列は残りの行インデックス $\tilde{S} = \{1, 3, 5\}$ から構成される。
アクティブ行列と潜在行列の混合ブロック $H_X (\tilde{H}_Z)^T$ を考えた場合、行インデックスは $i \in S'$、列インデックスは $k \in \tilde{S}$ となる。このとき、インデックス差分 $r = (k - i) \pmod 6$ が取り得る値の集合は、
\[
\{(k - i) \pmod 6 \mid i \in \{0, 2, 4\}, k \in \{1, 3, 5\}\} = \{1, 3, 5\}
\]
となる。すなわち、混合ブロックは $\Psi_1, \Psi_3, \Psi_5$ のみから構成される。
潜在部がアクティブ部と直交しない（$H_X (\tilde{H}_Z)^T \neq 0$ かつ $H_Z (\tilde{H}_X)^T \neq 0$）ためには、$\Psi_1, \Psi_3, \Psi_5$ のうち少なくとも一つが非零行列でなければならない。

さらに、潜在行列の任意の低重みな行の線形結合が容易にゼロ空間に入り込まない（潜在ベースの距離 $d^{(lat)}$ を最大化する）ためには、単に非零であるだけでなく、これらの混合ブロックが十分なランクを持つ必要がある。

\subsubsection{探索アルゴリズムにおける両立手順}
上記の両立条件を探索アルゴリズムに組み込む手順は以下の通りである。

\begin{enumerate}
    \item \textbf{可換・非可換条件の明示的分割}\\
    パラメータの探索にあたり、インデックス差分 $r \in \{0, 2, 4\}$ に該当するペアに対しては厳密な可換性制約（線形合同式による候補の絞り込み）を適用し、$\Psi_0 = \Psi_2 = \Psi_4 = 0$ を担保する。
    
    \item \textbf{非可換ペアの戦略的配置}\\
    インデックス差分 $r \in \{1, 3, 5\}$ に該当するペア群のなかに、意図的に非可換となるペアを配置する。例えば、$\Psi_1, \Psi_3, \Psi_5$ のすべてが非零行列となるように、各 $r$ に対して少なくとも1組以上の非可換ペア $(F_u, G_{r-u})$ を設定する。
    
    \item \textbf{混合ブロックのランク評価}\\
    バックトラッキング探索の各ステップにおいて候補パラメータが仮決定された際、構成される $\Psi_1, \Psi_3, \Psi_5$ のランクや、それらから構成されるブロック行列のカーネル（零空間）の最小重みを評価する。低重みのベクトルがカーネルに属してしまう場合（すなわち $d^{(lat)}$ が想定より小さくなる場合）は、非直交性条件が不十分であると判定し、その探索枝を棄却（枝刈り）する。
\end{enumerate}

このように、インデックス集合を $S'=\{0, 2, 4\}$ に分散させることで、直交性維持のための制約を最小限に抑えつつ、残された自由度 $\{1, 3, 5\}$ をすべて潜在部の非直交性強化に割り当てることが可能となる。これにより、理論上 $P$ をより小さく設定しても、目的の条件を満たす行列ペアを発見しやすくなる。

\section{可換性テーブルの動的再設計と非可換ペアの拡張手順}

本手法は、特定の非可換ペアに固定されていた従来の可換性テーブルを動的に再設計し、探索空間を拡張することで、より小さなブロックサイズ $P$ を達成するものである。非可換ペアを増やすことで代数的な自由度は高まるが、同時に内周（girth）の低下やトラッピングセットの増加という構造的なリスクを伴う。そのため、以下の具体的な手順に従い、探索空間の拡張とグラフ構造の評価を統合して実行する。

\begin{enumerate}
    \item \textbf{非可換ペアパターンの一般化と集合定義}\\
    特定の2組（例えば $(0, 3)$ と $(1, 2)$）のみを非可換とする固定的な設計を廃止し、非直交性が許容されるインデックス差分（例えば $r \in \{1, 3, 5\}$）に対応するすべてのペア群の中から、非可換とするペアの組み合わせを選択可能にする。3組以上のペアを非可換とするパターンや、これまで探索されていなかった非対称な組み合わせを許容することで、制約を緩和し、より小さな $P$ で成立しうる新たなパラメータ空間を構築する。

    \item \textbf{探索アルゴリズムにおける制約の動的切り替え}\\
    バックトラッキング探索の各定式化ステップにおいて、現在評価中のペア $(F_u, G_{r-u})$ が設定した非可換パターンに含まれるか否かを判定する。含まれない場合は、線形合同式 $Ax \equiv B \pmod P$ の解を用いて厳密に可換性を強制し候補を絞り込む。一方、含まれる場合は、当該合同式の解を候補から明示的に除外することで非可換性を保証する。これにより、任意のパターンの可換性テーブルに基づく候補生成をプログラム上でシームレスに実行する。

    \item \textbf{目的関数に基づく構造的欠陥の早期枝刈り}\\
    非可換ペアの拡張に伴う最大のリスクは、復号において有害となるトラッピングセットの発生と、それに起因する低重み論理演算子の漏出である。これを防ぐため、探索中に仮決定されたパラメータを用いて部分的なタナーグラフを構成し、サイクルの数やトラッピングセットの形成状況を動的にスコアリングする。これらの構造的欠陥が規定の閾値を超過した場合、その探索枝は致命的なエラーフロアを引き起こすリスクが高いと判断し、深層に達する前に即座に棄却（枝刈り）する。
\end{enumerate}
\subsection{可換性テーブルの動的再設計と非可換ペアの拡張手順}

本手法は、特定の非可換ペアに限定されていた従来の設計を脱却し、可換性テーブルを動的に再設計することで探索空間を拡張するアプローチである。具体的には、アクティブ行列の直交性維持に影響を与えないインデックス差分（例えば $r \in \{1, 3, 5\}$）に対応するペア候補の中から、3組以上の非可換ペアや非対称なパターンをパラメータ集合 $\Omega$ として自由に定義可能とする。

探索アルゴリズムのバックトラッキング過程では、この集合 $\Omega$ に基づいて制約充足の論理を動的に切り替える。評価対象のペアが $\Omega$ に含まれない場合は線形合同式 $Ax \equiv B \pmod P$ の解を用いて厳密に可換性を強制し、含まれる場合は当該合同式の解を候補から明示的に除外することで非可換性を保証する。

さらに、非可換ペアの拡張に伴う構造的リスク（潜在的な低重み論理演算子の漏出や、BP復号において有害となるトラッピングセットの発生）を抑制するため、探索の各ステップで仮決定されたパラメータから混合ブロックの行列や部分的なタナーグラフを構成し評価を行う。混合ブロックの零空間の最小重み評価により非直交性の担保を確認するとともに、事前に定義した有害な基本トラッピングセット（ETS）の部分グラフ同型判定を動的に実行する。これらの構造的欠陥が許容閾値を超過した探索枝は、最終ステップに到達する前であっても即座に棄却（枝刈り）する。

これらの一連の手順を統合することで、可換性の制約を安全に緩和しつつ、直交性、距離、内周（girth）といった符号の性能要件を満たす、より小さなブロックサイズ $P$ を持つ行列構成の効率的な探索が可能となる。
\end{document}